\section{Introduction}
\label{sec:introduction}

\PARstart{P}{hasor} measurement units (PMUs) reveal electromechanical transients through synchronized voltage and current phasors, yet only a fraction of the buses and branches in a transmission network are normally instrumented \cite{phadke2008synchronized,zhao2019dse}. A disturbance may therefore originate outside the measured region, alter hidden voltages, and propagate to the available devices. At the same time, a missing frame or corrupted channel can resemble a physical transient. A useful diagnosis must answer three distinct questions: what is the network state, which physical or measurement event is present, and where did each applicable event originate?

Strong prior methods solve important subsets of this problem. Dynamic estimators combine physical models with synchrophasors \cite{zhao2019dse,katanic2024recursive}; sparse-PMU classifiers identify events and faulted assets \cite{pandey2020realtime,yildiz2024few}; and recent learning systems reconstruct hidden trajectories before localization or estimate state and attack location from partial observations \cite{rimkus2025gnn,zhai2025partial}. These results make a broad ``joint inference'' novelty claim untenable. The unresolved contract considered here is narrower: causal inference in an under-observed transmission system when physical interventions and measurement-integrity failures can coexist, the source asset may be absent from source-labeled training, and the sensing geometry may make two candidate sources physically indistinguishable.

This distinction matters because event recognition, state reconstruction, and source localization have different information requirements. A generation change can be detected throughout the network while its source remains ambiguous. A dropped PMU identifies an integrity location directly but removes information needed to localize a simultaneous physical event. Likewise, a state estimate can have low point error although the nuisance intervention is not uniquely identifiable. Treating all three tasks as one multiclass label hides these differences and encourages scores whose denominator does not match the claimed capability.

We represent a physical event as an intervention hypothesis $h=(c,\ell,\tau,\gamma)$ containing family, source, onset, and magnitude, and represent measurement integrity separately through a time-varying availability operator and an additive corruption process. Each physical candidate produces a counterfactual DAE response. Candidate comparison profiles the admissible initial-state and integrity nuisance directions. The resulting separation margin exposes source pairs that the available PMUs cannot distinguish and motivates a candidate set or abstention instead of a forced source.

The empirical record currently contains two complementary, frozen campaigns. The PowerDynamics campaign audits an IEEE 39-bus hybrid DAE, reconstructs 31 hidden complex bus voltages from eight voltage/current PMUs, and studies model mismatch and physical recentering. The separate ANDES campaign contains \LegacyEventTrajectories{} trajectories spanning eight abnormal classes plus normal operation and evaluates a physics-guided hierarchical ExtraTrees baseline. The two campaigns share PMU buses and a network family, but not simulator semantics, state coordinates, or split contracts; their metrics are never pooled. The event study is presented as a baseline and failure analysis, not as evidence that the proposed joint posterior has already been executed.

The paper makes four contributions.
\begin{enumerate}
  \item We formulate sparse-PMU monitoring as joint inference over hidden state, physical intervention, physical source, and measurement integrity. The eight event labels are mapped to composable physical and integrity variables rather than treated as eight unrelated phenomena.
  \item We distinguish functional observability of the requested voltage field from event-source diagnosability. A finite-horizon affine-manifold distance gives a local theoretical limit: intersecting hypothesis manifolds admit nuisance configurations that no decision rule can distinguish uniformly.
  \item We consolidate two audited evidence lines. On the state campaign, a preregistered structural residual predicts bus-level reconstruction difficulty and measurement-only regularized physical recentering recovers operating-point-shifted voltages. On the event campaign, multiscale and spatial features improve known-source localization, but complete source holdout collapses to \HeldSourceCorrect{}/\HeldSourceDecisions{} exact decisions.
  \item We define the reproducible comparison required for the next joint experiment: learned-only diagnosis, physics-only hypothesis inference, physics plus learned residual discrepancy, and diagnosability-aware inference with abstention, all evaluated on source-disjoint parents with explicit alarm exposure and PMU loss/join conditions.
\end{enumerate}

The strongest result is also the main limitation. The state component is accurate under the tested nominal and operating-point-shift contracts, and the known-source event baseline attains detection F1 \LegacyEventDetectionFOne{} and event macro-F1 \LegacyEventMacroFOne{}. Yet source localization does not transfer to assets excluded from training. The current manuscript therefore establishes the mathematical and empirical boundary that a joint physics-informed method must cross; it does not convert a proposed end-to-end estimator into a completed result.

\begin{figure*}[t]
\centering
\resizebox{\textwidth}{!}{\begin{tikzpicture}[
  x=1cm,y=1cm,
  font=\sffamily\scriptsize,
  >=Latex,
  box/.style={draw=black!55,rounded corners=2pt,thick,align=center,
    minimum height=0.72cm,inner sep=3pt,fill=white},
  state/.style={box,draw=accessblue!75!black,fill=accessblue!8},
  event/.style={box,draw=orange!78!black,fill=orange!9},
  loadcase/.style={box,draw=green!55!black,fill=green!7},
  pending/.style={box,draw=purple!72!black,densely dashed,fill=purple!5},
  arr/.style={->,thick,draw=black!60,rounded corners=2pt},
  pendingarr/.style={->,thick,densely dashed,draw=purple!72!black,rounded corners=2pt},
  lane/.style={font=\sffamily\tiny\bfseries,anchor=east,align=right}
]

% Column headings make the contract readable before the individual boxes.
\foreach \x/\txt in {0.95/OBSERVATIONS,3.72/PHYSICAL MODEL,6.72/INFERENCE,9.82/EVIDENCE,13.10/OUTPUT}{
  \node[font=\sffamily\tiny\bfseries,text=black!58] at (\x,3.48) {\txt};
}

\node[state,text width=2.05cm,minimum height=1.15cm] (pmu) at (0.95,1.48)
  {PMU frame at $k$\\$\bm y_k=\{V,I,f\}$\\$\bm M_k,\bm b_k$, quality flags};

\node[state,text width=2.10cm] (dae) at (3.72,2.45)
  {Network and DAE model\\$\bm f,\bm g,\bm h_{\mathcal P}$};
\node[event,text width=2.10cm] (bank) at (3.72,0.45)
  {Physical hypothesis bank\\$h=(c,\ell,\tau,\gamma)$};

\node[state,text width=2.25cm] (statefit) at (6.72,2.45)
  {Causal state estimator\\frozen B0--B3 tests};
\node[pending,text width=2.38cm,minimum height=0.90cm] (jointfit) at (6.72,1.20)
  {Candidate-conditioned fit\\state and integrity nuisance};
\node[loadcase,text width=2.25cm] (loadfit) at (6.72,-0.05)
  {Load tangent\\and fitted noise};

\node[pending,text width=2.30cm,minimum height=0.90cm] (compare) at (9.82,1.20)
  {Profiled candidate evidence\\$D_{ij},\ p(h\mid\bm y_{0:k})$};
\node[loadcase,text width=2.22cm] (loadpost) at (9.82,-0.05)
  {Closed evaluated bank\\$H_0$ and 16 load buses};

\node[state,text width=2.32cm] (stateout) at (13.10,2.45)
  {Hidden-voltage estimate\\$\widehat{\bm z}_{\mathcal U,k},\bm\Sigma_{z,k}$};
\node[event,text width=2.32cm] (eventout) at (13.10,1.05)
  {Event, source, integrity\\$\widehat c_k,\widehat\ell_k,\widehat q_k$};
\node[pending,text width=1.55cm] (abstain) at (15.55,1.05)
  {candidate set\\or abstain};

% Three distinct lanes: evaluated state, pending joint inference, and evaluated load-only inference.
\node[lane,text=accessblue!80!black] at (-0.35,2.45) {evaluated\\state lane};
\node[lane,text=purple!78!black] at (-0.35,1.20) {joint lane\\not evaluated};
\node[lane,text=green!48!black] at (-0.35,-0.05) {evaluated\\load lane};

\begin{scope}[on background layer]
  \draw[arr,draw=accessblue!72!black] ($(pmu.east)+(0,0.40)$) -- (dae.west);
  \draw[arr,draw=accessblue!72!black] (dae) -- (statefit);
  \draw[arr,draw=accessblue!72!black] (statefit) -- (stateout);

  \draw[pendingarr] (pmu.east) -- (jointfit.west);
  \draw[pendingarr] (dae.south) -- (jointfit.north);
  \draw[pendingarr] (bank.north) -- (jointfit.south);
  \draw[pendingarr] (jointfit) -- (compare);
  \draw[pendingarr] (compare.east) -- (eventout.west);
  \draw[pendingarr] (compare.north east) -| (stateout.south);

  \draw[arr,draw=green!48!black] (pmu.south) -- (0.95,-0.62) -- (5.24,-0.62) |- (loadfit.west);
  \draw[arr,draw=green!48!black] (bank) -- (loadfit);
  \draw[arr,draw=green!48!black] (loadfit) -- (loadpost);
  \draw[arr,draw=green!48!black] (loadpost.east) -| (eventout.south);
  \draw[pendingarr] (eventout) -- (abstain);
\end{scope}

\node[font=\sffamily\tiny,text=black!58,anchor=west] at (1.55,-0.92)
  {solid blue/green: executed evidence};
\node[font=\sffamily\tiny,text=purple!78!black,anchor=west] at (6.45,-0.92)
  {purple dashed: proposed multi-family joint path};
\node[font=\sffamily\tiny,text=orange!80!black,anchor=west] at (12.55,-0.92)
  {orange: physical/integrity decision};
\end{tikzpicture}
}
\caption{Joint inference contract. The physical model generates candidate state and PMU responses, while availability and quality govern which residuals are admissible. Physical interventions and measurement-integrity failures remain separate latent variables. The rightmost outputs distinguish hidden-state uncertainty from event/source ambiguity and permit a candidate set instead of a forced source. Solid arrows describe the proposed online information flow; this paper validates its state-estimation and historical diagnosis components separately.}
\label{fig:joint-architecture}
\end{figure*}

The remainder of the paper reviews the closest estimation and localization contracts, defines the hybrid DAE and eight event labels, derives the observability and diagnosability limits, specifies the estimator hierarchy, documents both evidence campaigns, and reports the validated results and unresolved joint-inference gate.
